%! The Quadratic Formula & the Discriminant — Unit 2, Lesson 2.5 — Group Activity
\documentclass[10pt]{article}
\usepackage{algebra2-article}
\usepackage{algebra2-boxes}
\newcommand{\TallMath}[1]{$\displaystyle #1\rule[-1.4em]{0pt}{3.2em}$}
\newcommand{\lgrid}[4]{%
  \draw[step=1, linegray!40, very thin] (#1,#3) grid (#2,#4);
  \draw[->, semithick, charcoal] (#1,0)--(#2,0) node[below right, font=\scriptsize]{$x$};
  \draw[->, semithick, charcoal] (0,#3)--(0,#4) node[left, font=\scriptsize]{$y$};}

\begin{document}

\pageheader{Unit 2, Lesson 2.5}{Group Activity}

\vspace{0.08in}

\begin{headlinebox}{forestbg}
\small\textbf{One Formula, Every Quadratic.} With your partner, always start in \textbf{standard form}
$ax^2+bx+c=0$, read $a,b,c$ \emph{with signs}, and compute the \textbf{discriminant} $b^2-4ac$ \emph{first}
--- it tells you what kind of answer to expect. Then run the formula and simplify.
\end{headlinebox}

\vspace{0.06in}

\begin{tcolorbox}[colback=white, colframe=black!40, title=\textbf{Tier R --- Remediate},
    fonttitle=\bfseries, arc=1.5mm, left=3mm, right=3mm, top=2mm, bottom=2mm, breakable]
\textbf{1. Standard form, then $a,b,c$.}
\begin{itemize}\setlength{\itemsep}{3pt}
  \item $x^2-5x=-6 \Rightarrow$ \blank{2.8cm} \hfill $a=\blank{0.9cm}\ b=\blank{0.9cm}\ c=\blank{0.9cm}$
  \item $3x^2+2=7x \Rightarrow$ \blank{2.8cm} \hfill $a=\blank{0.9cm}\ b=\blank{0.9cm}\ c=\blank{0.9cm}$
\end{itemize}

\textbf{2. Compute the discriminant $b^2-4ac$.}
\[ x^2-4x+1=0:\ \blank{1.6cm} \qquad x^2+2x+5=0:\ \blank{1.6cm} \qquad x^2-6x+9=0:\ \blank{1.6cm} \]

\textbf{3. Classify the roots} from the discriminant (circle one for each).
\begin{center}
\renewcommand{\arraystretch}{1.5}
\begin{tabularx}{\linewidth}{@{}c X@{}}
\textbf{$b^2-4ac$} & \textbf{Roots (circle one)} \\
\hline
$25$  & two real \quad / \quad one repeated real \quad / \quad two complex \\
$0$   & two real \quad / \quad one repeated real \quad / \quad two complex \\
$-9$  & two real \quad / \quad one repeated real \quad / \quad two complex \\
\end{tabularx}
\end{center}

\textbf{4. Match to a graph.} Which parabola goes with $b^2-4ac>0$, $=0$, and $<0$?
\begin{center}
\begin{tikzpicture}[scale=0.34]
  \lgrid{-2}{2}{-2}{4}
  \begin{scope}\clip (-2,-2) rectangle (2,4);
    \draw[very thick, forest, domain=-1.7:1.7, samples=50, smooth] plot (\x,{\x*\x-1});
  \end{scope}
  \node[font=\scriptsize\bfseries, charcoal] at (0,-3.2){A};
\end{tikzpicture}
\hspace{0.6cm}
\begin{tikzpicture}[scale=0.34]
  \lgrid{-2}{2}{-1}{4}
  \begin{scope}\clip (-2,-1) rectangle (2,4);
    \draw[very thick, forest, domain=-1.6:1.6, samples=50, smooth] plot (\x,{\x*\x+1});
  \end{scope}
  \node[font=\scriptsize\bfseries, charcoal] at (0,-2.2){B};
\end{tikzpicture}
\hspace{0.6cm}
\begin{tikzpicture}[scale=0.34]
  \lgrid{-1}{3}{-1}{4}
  \begin{scope}\clip (-1,-1) rectangle (3,4);
    \draw[very thick, forest, domain=-0.6:2.6, samples=50, smooth] plot (\x,{(\x-1)*(\x-1)});
  \end{scope}
  \node[font=\scriptsize\bfseries, charcoal] at (1,-2.2){C};
\end{tikzpicture}
\end{center}
\[ b^2-4ac>0:\ \blank{1.2cm} \qquad b^2-4ac=0:\ \blank{1.2cm} \qquad b^2-4ac<0:\ \blank{1.2cm} \]
\end{tcolorbox}

\begin{tcolorbox}[colback=white, colframe=black!40, title=\textbf{Tier A --- Approaching Proficiency},
    fonttitle=\bfseries, arc=1.5mm, left=3mm, right=3mm, top=2mm, bottom=2mm, breakable]
\textbf{1. Solve with the formula.} Give simplest radical or $a+bi$ form.

\textbf{(a)} $x^2-3x-10=0$
\begin{work}
  x &= \frac{3\pm\sqrt{(-3)^2-4(1)(-10)}}{2(1)} \\
    &= \frac{3\pm\sqrt{49}}{2} \\
    &= \frac{3\pm7}{2} \\
    &= 5,\ -2
\end{work}
\textbf{(b)} $x^2-4x+1=0$
\begin{work}
  x &= \frac{4\pm\sqrt{(-4)^2-4(1)(1)}}{2(1)} \\
    &= \frac{4\pm\sqrt{12}}{2} \\
    &= \frac{4\pm2\sqrt{3}}{2} \\
    &= 2\pm\sqrt{3}
\end{work}

\textbf{2. Complex cases.} The discriminant is negative --- keep both conjugate roots.

\textbf{(a)} $x^2-2x+5=0$
\begin{work}
  x &= \frac{2\pm\sqrt{(-2)^2-4(1)(5)}}{2(1)} \\
    &= \frac{2\pm\sqrt{-16}}{2} \\
    &= \frac{2\pm4i}{2} \\
    &= 1\pm2i
\end{work}
\textbf{(b)} $2x^2+3x+2=0$
\begin{work}
  x &= \frac{-3\pm\sqrt{3^2-4(2)(2)}}{2(2)} \\
    &= \frac{-3\pm\sqrt{-7}}{4} \\
    &= \frac{-3\pm i\sqrt{7}}{4} \\
    &= -\frac{3}{4}\pm\frac{\sqrt{7}}{4}i
\end{work}

\textbf{3. Verify.} Substitute $x=5$ into $x^2-3x-10$ (from \#1) and show it gives $0$.
\begin{work}
  (5)^2-3(5)-10 &= 25-15-10 \\
                &= 0
\end{work}
\end{tcolorbox}

\begin{tcolorbox}[colback=white, colframe=black!40, title=\textbf{Tier E --- Extension},
    fonttitle=\bfseries, arc=1.5mm, left=3mm, right=3mm, top=2mm, bottom=2mm, breakable]
\textbf{1. Model (projectile).} A rocket's height is $h(t)=-16t^2+32t+20$ (feet, $t$ in seconds). Find when
it hits the ground ($h=0$) with the formula, and keep only the sensible root.
\begin{work}
  t &= \frac{-32\pm\sqrt{32^2-4(-16)(20)}}{2(-16)} \\
    &= \frac{-32\pm\sqrt{2304}}{-32} \\
    &= \frac{-32\pm48}{-32} \\
    &= -0.5,\ 2.5 \quad\text{reject } t=-0.5,\ \text{so it lands at } t=2.5\text{ s}
\end{work}

\tcbbreak
\textbf{2. Discriminant as a dial (parameter).}
\begin{itemize}\setlength{\itemsep}{3pt}
  \item For what value(s) of $k$ does $x^2+kx+9=0$ have exactly \emph{one repeated} real root? \blank{2.6cm}
  \item For what values of $k$ does $x^2-4x+k=0$ have \emph{two complex} roots? \blank{2.6cm}
\end{itemize}

\textbf{3. Verify a complex root.} Substitute $x=1+2i$ into $x^2-2x+5$ (Tier A \#2) and show it gives $0$.
\begin{work}
  (1+2i)^2-2(1+2i)+5 &= (1+4i+4i^2)-(2+4i)+5 \\
                     &= (1+4i-4)-(2+4i)+5 \\
                     &= (-3+4i)-(2+4i)+5 \\
                     &= 0
\end{work}
\end{tcolorbox}

\end{document}
